% !TEX root = ../../ctfp-print.tex

\lettrine[lhang=0.17]{圏}{は}恥ずかしいほど単純な概念だ。
圏は\newterm{対象}と、それらの間を結ぶ\newterm{射}（矢印）から構成される。だから
こそ、圏を図で表現するのはとても簡単だ。対象は円や点として描け、射は\ldots{}矢印
だ。（たまには変化をつけて、対象をブタ、射を花火として描くこともある。）しかし圏
の本質は\emph{合成}にある。あるいは、合成の本質が圏だと言ってもいい。射は合成
できる。つまり、対象$A$から対象$B$への射があり、$B$から$C$への射があれば、
$A$から$C$へ向かう射――それらの合成――が必ずなければならない。

\begin{figure}
  \centering
  \includegraphics[width=0.8\textwidth]{images/img_1330.jpg}
  \caption{圏において、$A$から$B$への射と$B$から$C$への射があれば、
    $A$から$C$への直接の射（それらの合成）も存在しなければならない。
    この図は恒等射が欠けているため（後述）、完全な圏ではない。}
\end{figure}

\section{射としての関数}

もう抽象的すぎると感じているだろうか。絶望しないでほしい。具体的な話をしよう。
矢印――\newterm{射}とも呼ばれる――を関数として考えてみよう。型$A$の引数を取り
$B$を返す関数$f$がある。$B$を取り$C$を返す関数$g$もある。$f$の結果を$g$に渡す
ことでそれらを合成できる。これで$A$を取り$C$を返す新しい関数を定義したことになる。

数学では、このような合成は関数の間に小さな円を挟んで表記する：$g \circ f$。合成
の順序が右から左になっていることに注意してほしい。混乱を感じる人もいるかもしれ
ない。Unixのパイプ記法に慣れている人もいるだろう：

\begin{snip}{text}
lsof | grep Chrome
\end{snip}
あるいはF\#の山形記号\code{>>}も同様だ。これらはどちらも左から右へと進む。しかし
数学やHaskellでは、関数は右から左に合成される。$g \circ f$を「$f$の後の$g$」と
読むと理解しやすい。

Cコードを書いてより明確にしよう。型\code{A}の引数を取り型\code{B}の値を返す
関数\code{f}がある：

\begin{snip}{text}
B f(A a);
\end{snip}
そしてもう一つ：

\begin{snip}{text}
C g(B b);
\end{snip}
それらの合成は：

\begin{snip}{text}
C g_after_f(A a)
{
    return g(f(a));
}
\end{snip}
ここでも右から左への合成が見られる：\code{g(f(a))}、今度はCで。

2つの関数を取りその合成を返すテンプレートがC++標準ライブラリにあればいいのだが、
残念ながら存在しない。気分転換にHaskellを試してみよう。$A$から$B$への関数の
宣言はこうだ：

\src{snippet01}
同様に：

\src{snippet02}
それらの合成は：

\src{snippet03}
Haskellでいかに物事が単純かを見ると、C++で素直な関数的概念を表現できないのが
少し恥ずかしくなる。実際、HaskellではUnicode文字を使えるので、合成をこう書ける：
% don't 'mathify' this block
\begin{snip}{text}
g ◦ f
\end{snip}

Unicodeのダブルコロンや矢印も使える：
% don't 'mathify' this block
\begin{snipv}
f \ensuremath{\Colon} A → B
\end{snipv}
これが最初のHaskellのレッスンだ：ダブルコロンは「〜の型を持つ」という意味だ。
関数の型は2つの型の間に矢印を挿入することで作られる。2つの関数を合成するには
それらの間にピリオドを挿入する（またはUnicodeの円）。

\section{合成の性質}

任意の圏における合成が満たさなければならない2つの非常に重要な性質がある。

\begin{enumerate}
  \item
        合成は結合的である。合成可能な（つまり対象が端から端まで一致する）3つの射
        $f$、$g$、$h$があるとき、それらを合成するのに括弧は不要だ。数学的記法では
        こう表現される：
        \[h \circ (g \circ f) = (h \circ g) \circ f = h \circ g \circ f\]
        （疑似）Haskellで書くと：

        \src{snippet04}[b]
        （「疑似」と言ったのは、関数に対して等号が定義されていないからだ。）

        関数を扱う場合、結合性はかなり明らかだが、他の圏では明白でないかもしれない。

  \item
        すべての対象$A$に対して、合成の単位となる射が存在する。この射は対象からそれ
        自身へとループする。合成の単位であるということは、$A$で始まるまたは$A$で
        終わる任意の射と合成したとき、元の射をそのまま返すことを意味する。対象$A$の
        単位射は$\idarrow[A]$（$A$上の\newterm{恒等射}）と呼ばれる。$f$が$A$から
        $B$へ向かうとき、数学的記法では：
        \[f \circ \idarrow[A] = f\]
        そして
        \[\idarrow[B] \circ f = f\]
\end{enumerate}
関数を扱う場合、恒等射は引数をそのまま返す恒等関数として実装される。この実装は
すべての型で同じだ。つまり、この関数は全称多相的である。C++ではテンプレートとして
定義できる：

\begin{snip}{cpp}
template<typename T> T id(T x) { return x; }
\end{snip}
もちろん、C++では物事はそれほど単純ではない。渡すものだけでなく、どのように渡す
か（値渡し、参照渡し、const参照渡し、ムーブなど）も考慮しなければならないからだ。

Haskellでは、恒等関数は標準ライブラリ（Preludeと呼ばれる）の一部だ。その宣言と
定義はこうなる：

\src{snippet05}
見てわかるように、Haskellの多相関数は朝飯前だ。宣言では型を型変数に置き換えるだ
けでいい。コツはこうだ：具体的な型の名前は常に大文字で始まり、型変数の名前は
小文字で始まる。したがってここでの\code{a}はすべての型を表す。

Haskellの関数定義は、関数名に続いて仮引数――ここでは\code{x}の1つだけ――が来る。
関数の本体は等号の後に続く。この簡潔さは初心者をしばしば驚かせるが、すぐに完全に
理にかなっていることがわかるだろう。関数定義と関数呼び出しは関数型プログラミングの
基本中の基本なので、その構文は最小限に削ぎ落とされている。引数リストの周りに括弧
がないだけでなく、引数の間にカンマもない（複数の引数を持つ関数を定義するときに
後ほど見ることになる）。

関数の本体は常に式だ――関数に文はない。関数の結果はこの式だ――ここでは単に
\code{x}。

これで2番目のHaskellのレッスンを終える。

恒等条件は（再び疑似Haskellで）こう書ける：

\src{snippet06}
こう自問するかもしれない：なぜ恒等関数――何もしない関数――を気にする必要があるのか？
同様に、なぜ私たちはゼロという数を気にするのか？ゼロは「何もない」ことのシンボル
だ。古代ローマ人はゼロのない数体系を持っていたが、今日まで残っているものを含む、
優れた道路や水道橋を建設できた。

ゼロや$\id$のような中立的な値は、シンボリック変数を扱うときに非常に役立つ。だから
ローマ人は代数があまり得意でなかったが、ゼロの概念に精通していたアラブ人やペルシャ
人は優れていた。だから恒等関数は、高階関数への引数として、あるいは高階関数からの
戻り値として、非常に便利になる。高階関数こそが関数のシンボリックな操作を可能にする
ものだ。それらは関数の代数だ。

まとめると：圏は対象と射（矢印）から構成される。射は合成でき、合成は結合的だ。
すべての対象は合成の単位として機能する恒等射を持つ。

\section{合成はプログラミングの本質}

関数型プログラマは問題へのアプローチに独特の方法がある。非常に禅的な問いから
始めるのだ。例えば、インタラクティブなプログラムを設計するとき、「インタラクション
とは何か？」と問うだろう。コンウェイのライフゲームを実装するとき、おそらく生の
意味について熟考するだろう。この精神で、私はこう問いたい：プログラミングとは何か？

最も基本的なレベルでは、プログラミングはコンピュータに何をすべきか伝えることだ。
「メモリアドレスxの内容を取り、それをレジスタEAXの内容に加算せよ」。しかしアセン
ブリでプログラムするときでも、コンピュータに与える命令はより意味のある何かの表現
だ。私たちは非自明な問題を解決している（自明なら、コンピュータの助けは不要だ）。
そして問題をどのように解決するか？より大きな問題をより小さな問題に分解する。小さな
問題がまだ大きすぎるなら、さらに分解する、というように。最終的に、すべての小さな
問題を解くコードを書く。そしてプログラミングの本質がくる：より大きな問題への解決策
を作るために、それらのコードの断片を合成するのだ。断片を元に戻せなければ、分解は
意味をなさない。

この階層的な分解と再合成のプロセスは、コンピュータによって強いられるものではない。
それは人間の心の限界を反映している。私たちの脳は一度に少数の概念しか扱えない。
心理学で最も引用される論文の一つである
\urlref{http://en.wikipedia.org/wiki/The_Magical_Number_Seven,_Plus_or_Minus_Two}{マジカル
  ナンバー7±2}は、私たちが心の中に保持できる「チャンク」の情報は$7 \pm 2$個だけ
だと仮定した。人間の短期記憶の理解の詳細は変わっているかもしれないが、それが限られ
ていることは確かだ。要するに、私たちはオブジェクトのごたまぜやコードのスパゲッティ
を扱えない。よく構造化されたプログラムが見た目に美しいからではなく、そうしなければ
私たちの脳が効率的に処理できないから、構造が必要なのだ。コードの一部を「エレガント」
や「美しい」と表現することがよくあるが、実際に意味しているのは、それが私たちの限ら
れた人間の心で処理しやすいということだ。エレガントなコードは、私たちの精神的な消化
システムが同化するのにちょうどよいサイズとちょうどよい数のチャンクを作り出す。

では、プログラムの合成にとって適切なチャンクとは何か？それらの表面積は体積よりも
ゆっくりと増加しなければならない。（幾何学的オブジェクトの表面積はサイズの二乗で
増加し、体積はサイズの三乗で増加するという直感があるので、私はこのアナロジーが
好きだ。）表面積はチャンクを合成するために必要な情報だ。体積はそれらを実装するため
に必要な情報だ。アイデアはこうだ：チャンクが一度実装されれば、その実装の詳細を忘れ、
他のチャンクとどのように相互作用するかに集中できる。オブジェクト指向プログラミング
では、表面はオブジェクトのクラス宣言または抽象インターフェースだ。関数型プログラミ
ングでは、関数の宣言だ。（少し単純化しているが、それが要点だ。）

圏論は、対象の内部を覗くことを積極的に思いとどまらせるという意味で極端だ。圏論に
おける対象は抽象的で捉えどころのない実体だ。その対象について知ることができるのは、
他の対象とどのように関係しているか――射を使ってどのように繋がっているか――だけだ。
これはインターネット検索エンジンが受信リンクと送信リンクを分析してウェブサイトを
ランク付けする方法だ（不正をしないとき）。オブジェクト指向プログラミングでは、
理想的なオブジェクトは抽象インターフェース（純粋な表面、体積なし）を通じてのみ
見え、メソッドが射の役割を果たす。他のオブジェクトとどのように合成するかを理解
するためにオブジェクトの実装を掘り下げなければならなくなった瞬間、プログラミング
パラダイムの利点を失ったことになる。

\section{課題}

\begin{enumerate}
  \tightlist
  \item
        好きな言語（もし好きな言語がHaskellであれば、2番目に好きな言語）で、
        できる限り恒等関数を実装せよ。
  \item
        好きな言語で合成関数を実装せよ。2つの関数を引数に取り、それらの合成で
        ある関数を返す。
  \item
        合成関数が恒等性を尊重することをテストしようとするプログラムを書け。
  \item
        ワールドワイドウェブはある意味で圏だろうか？リンクは射だろうか？
  \item
        Facebookは、人々を対象、友情を射とする圏だろうか？
  \item
        有向グラフはいつ圏になるか？
\end{enumerate}
